\subsection{\texorpdfstring{第~\ref{sec:hopf-rinow-theorem}~节}{Hopf-Rinow 定理}}\label{subsec:solutions-hopf-rinow-theorem}

\begin{description}[font=\normalfont,leftmargin=2em,labelwidth=1.5em]
\RGsolution{ex:proof-of-riemannian-metric-space-topology}{riemannian-metric-space-topology}
为证明 Proposition~\ref{prop:riemannian-distance-defines-metric-topology}, 先验证 $d$ 是一个距离:
\begin{enumerate}[label=(\roman*),ref=\roman*,leftmargin=*]
\item 因为 $d(p,q)$ 是一组非负数的下确界, 所以显然有 $d(p,q)\ge0$. 也容易验证 $d(p,p)=0$. 设 $B_\varepsilon(p)$ 为正规球, 并设 $q\ne p$. 若 $q\notin B_\varepsilon(p)$, 则 $d(p,q)\ge\varepsilon>0$. 若 $q\in B_\varepsilon(p)$, 则 $d(p,q)=\lVert\exp_p^{-1}(q)\rVert>0$. 因此 $d(p,q)=0$ 当且仅当 $p=q$.
\item 若 $\gamma:[a,b]\to M$ 是连接 $p$ 与 $q$ 的分段可微曲线, 则由 $\widetilde\gamma(t)=\gamma(a+b-t)$ 定义的 $\widetilde\gamma:[a,b]\to M$ 是连接 $q$ 与 $p$ 的分段可微曲线, 且长度相同. 因此 $d(p,q)=d(q,p)$.
\item 对每一对分别连接 $p$ 与 $q$, 以及 $q$ 与 $r$ 的分段可微曲线 $\gamma_1$ 和 $\gamma_2$, 容易构造连接 $p$ 与 $r$ 的分段可微曲线 $\gamma$, 使 $l(\gamma)=l(\gamma_1)+l(\gamma_2)$. 取下确界, 得 $d(p,r)\le d(p,q)+d(q,r)$.
\end{enumerate}
为证明 $M$ 上由度量诱导的拓扑是通常的拓扑, 注意正规球构成度量空间拓扑的一组基, 且这些球都是通常拓扑中的开集. 另一方面, 设 $U\subset M$ 是通常拓扑中的开集, 且 $p\in U$. 若 $B_\varepsilon(p)$ 是正规球, 则 $\exp_p^{-1}(U\cap B_\varepsilon(p))$ 是 $T_pM$ 中含原点的开集, 因而对某个 $\delta>0$, 它包含 $B_\delta(0)$, 这蕴含 $B_\delta(p)\subset U$. 这说明正规球也是通常拓扑的一组基, 所以两个拓扑一致.

\RGsolution{ex:metric-ball-in-slit-plane}{metric-ball-in-slit-plane}
集合 $B_7(0,4)$ 如 Fig.~\ref{fig:metric-ball-in-slit-plane} 所示. 较大的圆半径为 7, 较小的圆半径为 2.
\RGfigure[.56\textwidth]{metric-ball-in-slit-plane}{$\mathbb{R}^2\setminus\{(x,0)\mid-3\le x\le3\}$ 上的 $B_7(0,4)$}

\RGsolution{ex:ambrose-and-cartan-hadamard-theorems}{ambrose-and-cartan-hadamard-theorems}
\begin{enumerate}[label=(\alph*),ref=\alph*,leftmargin=*]
\item\label{sol:ambrose-local-isometry-covering} 因为 $f$ 是局部等距映射, 所以它是局部微分同胚. 为证明它是覆盖映射, 只需证明其为满射, 且对每个 $q\in N$, 都存在邻域 $U\ni q$, 使得 $f^{-1}(U)$ 是一些经 $f$ 与 $U$ 微分同胚的开集的不交并. 因为 $(M,g)$ 完备且 $f$ 将测地线送到测地线, 所以 $(N,h)$ 显然完备且 $f$ 为满射. 取任意 $q\in N$, 选取正规球 $U:=B_\varepsilon(q)$ (取充分小的 $\varepsilon>0$). 则
\[
  f^{-1}(U)=\bigcup_{p\in f^{-1}(q)}B_\varepsilon(p),
\]
而 $f$ 限制在 $B_\varepsilon(p)$ 上显然是微分同胚 (因为 $f$ 将过 $p$ 的测地线送到过 $q$ 的测地线, 故它在正规坐标下就是恒等映射). 最后, 若 $p_1,p_2\in f^{-1}(q)$ 且 $p_1\ne p_2$, 则 $B_\varepsilon(p_1)\cap B_\varepsilon(p_2)=\varnothing$, 因为交集中的点与 $p_1$ 和 $p_2$ 的距离必相等 (所以不能构成一个开集).
\item 由 $f(t)=e^{it}$ 定义的 $f:(0,+\infty)\to S^1$ 是满的局部等距映射 (对于通常的度量), 但它不是覆盖映射: 对点 $1\in S^1$, 不存在邻域 $U$ 使得 $f^{-1}(U)$ 是一些经 $f$ 与 $U$ 微分同胚的开集的不交并.
\item 若 $(M,g)$ 的曲率非正, 则由 Exercise~\ref{ex:minimizing-properties-of-geodesics}(\ref{ex:jacobi-fields-and-cut-locus}), 指数映射 $\exp_p:T_pM\to M$ 没有临界点, 即它是局部微分同胚. 在 $T_pM$ 上考虑 Riemann 度量 $h:=\exp_p^*g$. 这样选取使 $\exp_p$ 成为局部等距映射, 并将过原点的直线送到过 $p$ 的测地线. 由于这些测地线对参数的一切取值都有定义, 可知 $\exp_0:T_0T_pM\to T_pM$ 对所有 $v\in T_0T_pM$ 都有定义, 所以 $(T_pM,h)$ 完备. 由 (\ref{sol:ambrose-local-isometry-covering}), $\exp_p$ 是覆盖映射.
\end{enumerate}
\end{description}
